\lecture

\begin{guiding}
What is a ``continuously varying family of vector spaces'' over a space, how do
we prove that such a family is genuinely twisted, and which constructions of
linear algebra survive the passage from single vector spaces to families?
\end{guiding}

\section{Definitions and first examples}

We begin by fixing the stage. A \emph{topological manifold} is a Hausdorff space
$X$ in which every point has a neighbourhood homeomorphic to $\R^n$; it is
\emph{smooth} when the transition functions between charts are smooth. For a
smooth manifold the tangent space $T_xX$ is obtained from paths $\gamma$ through
$x$, with $\gamma(0)=x$, by taking velocities
\[
\frac{d}{dt}\Big|_{t=0}\gamma(t).
\]

The total space of all tangent spaces is the first and most important example of
the structure we are about to define.

\begin{proposition}
For a smooth $n$-manifold $X$, the set $TX=\bigsqcup_{x\in X}T_xX$ is a smooth
$2n$-manifold, the projection $\pi\colon TX\to X$ is smooth, each fibre
$\pi^{-1}(x)=T_xX$ is a real vector space, and each coordinate chart
$U\subset X$ gives a homeomorphism $U\times\R^n\to\pi^{-1}(U)$ which is linear
on each fibre.
\end{proposition}

\begin{definition}[Vector bundle]
A \emph{real vector bundle} $\xi$ over a topological space $B$ consists of

\begin{enumerate}
\item a topological space $E=E(\xi)$, the \emph{total space};
\item a continuous surjection $\pi\colon E\to B$, the \emph{projection};
\item a real vector space structure on every fibre $F_b(\xi):=\pi^{-1}(b)$,
\end{enumerate}

subject to the following \emph{local triviality} condition: every $b\in B$ has a
neighbourhood $U$, an integer $n\ge0$, and a homeomorphism
\[
h\colon U\times\R^n\longrightarrow \pi^{-1}(U)
\]
such that $x\mapsto h(b',x)$ is a vector space isomorphism
$\R^n\to F_{b'}(\xi)$ for every $b'\in U$. Such an $h$ is a \emph{local
trivialization}. If one may take $U=B$ the bundle is \emph{trivial}; if $n$ is
the same for all $b$ we speak of an \emph{$\R^n$-bundle}, or a bundle of
\emph{rank} $n$.
\end{definition}

\begin{example}\label{ex:mobius}
Three examples to keep in mind throughout.

\begin{enumerate}
\item The cylinder $S^1\times\R$ is the trivial line bundle over the circle.

\item The \emph{M\"obius bundle}: take $[0,1]\times\R$ and glue
$(0,t)\sim(1,-t)$. Locally it is a cylinder; globally it is twisted.

\item $TS^2$, an $\R^2$-bundle over $S^2$, which we will prove is non-trivial.
\end{enumerate}

\begin{center}
\begin{tikzpicture}[scale=0.95]
\draw[thick] (0,0) rectangle (2.4,1.4);
\draw[-{Stealth},thick] (0,0.5) -- (0,0.9);
\draw[-{Stealth},thick] (2.4,0.5) -- (2.4,0.9);
\node at (1.2,-0.55) {\small $S^1\times\R$: edges glued in the same direction};
\begin{scope}[xshift=6.4cm]
\draw[thick] (0,0) rectangle (2.4,1.4);
\draw[-{Stealth},thick] (0,0.5) -- (0,0.9);
\draw[-{Stealth},thick] (2.4,0.9) -- (2.4,0.5);
\node at (1.2,-0.55) {\small M\"obius: $(0,t)\sim(1,-t)$};
\end{scope}
\end{tikzpicture}
\end{center}
\end{example}

\begin{definition}
Two vector bundles $\xi,\eta$ over $B$ are \emph{isomorphic}, written
$\xi\iso\eta$, if there is a homeomorphism $f\colon E(\xi)\to E(\eta)$ commuting
with the projections and restricting to a linear isomorphism
$F_b(\xi)\to F_b(\eta)$ on every fibre.
\end{definition}

\begin{definition}
For a smooth submanifold $M\subset\R^n$, the \emph{normal bundle} $\nu$ of $M$
is the subspace of $M\times\R^n$ consisting of the pairs $(x,v)$ with
$v\perp T_xM$.
\end{definition}

\begin{definition}
A \emph{(cross-)section} of $\xi$ is a continuous map $s\colon B\to E(\xi)$ with
$s(b)\in F_b(\xi)$ for every $b$. It is \emph{nowhere zero} if $s(b)\neq0$ for
all $b$.
\end{definition}

Sections are our main probe. Note three things: a section of a tangent bundle is
exactly a vector field; every bundle has a zero section; and a trivial line
bundle has a nowhere zero section, obtained by transporting the constant section
$b\mapsto(b,1)$ through a trivialization.

\section{Sections and non-triviality}

\begin{guiding}
How can one \emph{prove} that a bundle is non-trivial? The strategy that will
serve us for the rest of the course: exhibit something that a trivial bundle
must possess, and show the bundle at hand does not possess it.
\end{guiding}

\begin{example}[The canonical line bundle]
Write $\RP^n=\{\{\pm x\}\mid x\in S^n\}$ and let $q\colon S^n\to\RP^n$ be the
quotient map. The \emph{canonical line bundle} $\gaman\to\RP^n$ is the subspace
of $\RP^n\times\R^{n+1}$ consisting of pairs $(\{\pm x\},v)$ with $v$ on the
line through $x$ and $-x$, that is, $v$ a multiple of $x$.
\end{example}

\begin{theorem}\label{thm:gamma-nontrivial}
For every $n\ge1$ the bundle $\gaman$ is not trivial.
\end{theorem}

\begin{proof}
\begin{strategy}
A trivial line bundle has a nowhere zero section, so it is enough to prove that
$\gaman$ has none. Pull a hypothetical section back to the sphere, where the
bundle becomes trivial: the section becomes an honest real-valued function,
antisymmetric under the antipodal map, and the intermediate value theorem
produces a zero.
\end{strategy}

\step{1} \emph{The pull-back of $\gaman$ to $S^n$ is trivial.} Consider
\[
q^*\gaman
=\big\{(x,v)\in S^n\times\R^{n+1}\ \big|\ v\in\R x\big\},
\]
\[
\Phi\colon q^*\gaman\longrightarrow S^n\times\R,
\qquad (x,tx)\longmapsto (x,t).
\]
The map $\Phi$ is a bijection, linear on fibres, and both it and its inverse
$(x,t)\mapsto(x,tx)$ are continuous, the former because $t=x\cdot v$ is a
polynomial in the coordinates.

\step{2} \emph{A section becomes a function.} Suppose $s\colon\RP^n\to E(\gaman)$
is a section. Then $x\mapsto s(q(x))$ is a section of $q^*\gaman$, so by Step 1
there is a continuous $t\colon S^n\to\R$ with
\[
s\big(q(x)\big)=\big(\{\pm x\},\ t(x)\,x\big)
\qquad\text{for all }x\in S^n.
\]

\step{3} \emph{The function is odd.} Since $q(-x)=q(x)$, evaluating the formula
of Step 2 at $-x$ gives $t(-x)\cdot(-x)=t(x)\cdot x$, whence
\[
t(-x)=-t(x).
\]

\step{4} \emph{Conclusion.} Fix $x\in S^n$ and let $\alpha$ be a path in $S^n$
from $x$ to $-x$, which exists because $S^n$ is path connected for $n\ge1$. The
continuous function $t\circ\alpha$ takes the values $t(x)$ and $-t(x)$ at its
endpoints. If $t(x)=0$ the section already vanishes at $q(x)$; otherwise the
two values have opposite signs and the intermediate value theorem gives
$x_0$ on the path with $t(x_0)=0$, so $s(q(x_0))=0$. In either case $s$ is not
nowhere zero, and $\gaman$ is not trivial.
\end{proof}

\begin{remark}
For $n=1$ the total space $E(\gamma^1_1)$ is precisely the M\"obius band of
Example~\ref{ex:mobius}: the theorem recovers the familiar fact that the
M\"obius band is twisted.
\end{remark}

The next lemma is the workhorse that lets us verify isomorphisms fibre by fibre,
without ever checking continuity of an inverse by hand.

\begin{lemma}\label{lem:fibrewise-iso}
Let $\xi,\eta$ be vector bundles over $B$ and let $f\colon E(\xi)\to E(\eta)$ be
continuous, commuting with the projections, and restricting to a linear
isomorphism $F_b(\xi)\to F_b(\eta)$ for every $b$. Then $f$ is a homeomorphism
and $\xi\iso\eta$.
\end{lemma}

\begin{proof}
\begin{strategy}
Bijectivity is immediate; the whole content is continuity of $f^{-1}$, which is
a local question. In local trivializations $f$ becomes multiplication by an
invertible matrix depending continuously on the base point, and Cramer's rule
inverts it continuously.
\end{strategy}

\step{1} \emph{$f$ is a bijection.} It is injective and surjective on each
fibre by hypothesis, and it preserves fibres, so it is a bijection on total
spaces.

\step{2} \emph{Reduction to a local statement.} Continuity of $f^{-1}$ may be
checked on the members of an open cover of $E(\eta)$. Fix $b\in B$ and choose a
neighbourhood $U$ over which both $\xi$ and $\eta$ are trivial, with
trivializations
\[
g\colon U\times\R^n\to\pi_\xi^{-1}(U),
\qquad
h\colon U\times\R^n\to\pi_\eta^{-1}(U).
\]
It suffices to prove that $f^{-1}$ is continuous on $\pi_\eta^{-1}(U)$.

\step{3} \emph{$f$ in coordinates.} The composite
$h^{-1}\circ f\circ g\colon U\times\R^n\to U\times\R^n$ preserves the first
coordinate and is linear in the second, so it has the form
\[
(b',x)\longmapsto\big(b',\,M(b')\,x\big),
\qquad M(b')\in\GL_n(\R).
\]
Each column $M(b')e_i$ is the second component of
$h^{-1}fg(b',e_i)$, hence depends continuously on $b'$; so all entries of
$M(b')$ are continuous functions of $b'$.

\step{4} \emph{Inverting continuously.} By Cramer's rule,
\[
M(b')^{-1}=\frac{1}{\det M(b')}\,\mathrm{adj}\,M(b'),
\]
whose entries are rational functions of the entries of $M(b')$ with non-vanishing
denominator. Hence $b'\mapsto M(b')^{-1}$ is continuous, and
$(b',y)\mapsto(b',M(b')^{-1}y)$ is continuous.

\step{5} \emph{Conclusion.} On $\pi_\eta^{-1}(U)$ we have
$f^{-1}=g\circ\big(h^{-1}fg\big)^{-1}\circ h^{-1}$, a composite of continuous
maps. Therefore $f^{-1}$ is continuous, $f$ is a homeomorphism, and it is a
fibrewise linear isomorphism by hypothesis.
\end{proof}

\begin{definition}
Sections $s_1,\dots,s_n$ of $\xi$ are \emph{nowhere dependent} if the vectors
$s_1(b),\dots,s_n(b)$ are linearly independent for every $b\in B$.
\end{definition}

\begin{theorem}\label{thm:trivial-iff-sections}
An $\R^n$-bundle $\xi$ is trivial if and only if it admits $n$ nowhere dependent
sections.
\end{theorem}

\begin{proof}
\begin{strategy}
In both directions the map relating $B\times\R^n$ and $E$ is the obvious one;
Lemma~\ref{lem:fibrewise-iso} does the work of checking it is an isomorphism.
\end{strategy}

\step{1} \emph{From sections to a trivialization.} Let $s_1,\dots,s_n$ be
nowhere dependent and define
\[
f\colon B\times\R^n\longrightarrow E(\xi),
\qquad
f(b,x)=x_1s_1(b)+\dots+x_ns_n(b).
\]
This is continuous, since addition and scalar multiplication are continuous on
$E$ and each $s_i$ is continuous; it commutes with the projections; and on the
fibre over $b$ it sends the standard basis to the basis $s_1(b),\dots,s_n(b)$ of
$F_b(\xi)$, hence is a linear isomorphism there. By
Lemma~\ref{lem:fibrewise-iso} it is an isomorphism of bundles, so $\xi$ is
trivial.

\step{2} \emph{From a trivialization to sections.} If
$h\colon B\times\R^n\to E(\xi)$ is a trivialization, set $s_i(b)=h(b,e_i)$. Each
$s_i$ is continuous, and for every $b$ the vectors $s_1(b),\dots,s_n(b)$ are the
image of a basis under a linear isomorphism, hence linearly independent.
\end{proof}

\begin{corollary}
$TS^2$ is non-trivial.
\end{corollary}

\begin{proof}
By Theorem~\ref{thm:trivial-iff-sections} a trivialization would provide two
nowhere dependent vector fields on $S^2$, in particular one nowhere zero vector
field, contradicting the hairy ball theorem.
\end{proof}

\section{Constructions}

\begin{guiding}
Which operations of linear algebra --- restriction, pull-back, products, direct
sums, subspaces, orthogonal complements --- make sense for whole families of
vector spaces? Essentially all of them, and this flexibility is what makes the
theory usable.
\end{guiding}

\begin{construction}[Restriction]
If $\xi$ is a bundle over $B$ and $\bar B\subset B$, then $\xi|_{\bar B}$, with
total space $\pi^{-1}(\bar B)$, is a vector bundle over $\bar B$: local
trivializations restrict. For smooth $\bar B\subset B$ this produces
$TB|_{\bar B}$.
\end{construction}

\begin{construction}[Pull-back]\label{con:pullback}
Given $\xi$ over $B$ and a continuous map $f\colon B_1\to B$, the
\emph{induced}, or \emph{pull-back}, bundle $f^*\xi$ over $B_1$ has total space
\[
E_1(f^*\xi)=\big\{(b,e)\in B_1\times E(\xi)\ \big|\ f(b)=\pi(e)\big\},
\]
topologized as a subspace, with projection $(b,e)\mapsto b$ and the vector space
structure of $F_{f(b)}(\xi)$ on the fibre over $b$. It fits into the square
\[
\begin{tikzcd}[row sep=large, column sep=large]
E_1(f^*\xi) \arrow[r,"\hat f"] \arrow[d] & E(\xi) \arrow[d,"\pi"] \\
B_1 \arrow[r,"f"] & B
\end{tikzcd}
\qquad
\hat f(b,e)=e .
\]
\end{construction}

\begin{proposition}
The pull-back $f^*\xi$ is a vector bundle.
\end{proposition}

\begin{proof}
\begin{strategy}
Only local triviality needs an argument, and a trivialization of $\xi$ near
$f(b)$ can simply be pulled back over the preimage.
\end{strategy}

\step{1} Fix $b\in B_1$ and choose a local trivialization
$h\colon U\times\R^n\to\pi^{-1}(U)$ of $\xi$ with $f(b)\in U$. Put
$U_1=f^{-1}(U)$, an open neighbourhood of $b$.

\step{2} Define
\[
h_1\colon U_1\times\R^n\longrightarrow \pi_1^{-1}(U_1),
\qquad
h_1(b_1,x)=\big(b_1,\ h(f(b_1),x)\big).
\]
This lands in $E_1$ because $\pi(h(f(b_1),x))=f(b_1)$, and it is continuous
because $f$ and $h$ are.

\step{3} It is bijective with continuous inverse
$(b_1,e)\mapsto(b_1,\ \mathrm{pr}_2\,h^{-1}(e))$, and for fixed $b_1$ it is the
linear isomorphism $x\mapsto h(f(b_1),x)$ onto $F_{f(b_1)}(\xi)$. Hence $h_1$ is
a local trivialization.
\end{proof}

\begin{definition}
A \emph{bundle map} from $\eta$ to $\xi$ is a continuous map
$g\colon E(\eta)\to E(\xi)$ carrying each fibre $F_b(\eta)$ isomorphically onto
a fibre $F_{b'}(\xi)$. It induces a map $\bar g\colon B(\eta)\to B(\xi)$,
$b\mapsto b'$, on base spaces, which is continuous.
\end{definition}

\begin{lemma}\label{lem:bundle-map-pullback}
If $g\colon E(\eta)\to E(\xi)$ is a bundle map with base map $\bar g$, then
$\eta\iso\bar g^{\,*}\xi$.
\end{lemma}

\begin{proof}
\begin{strategy}
There is only one reasonable map to write down; check its fibrewise behaviour
and appeal to Lemma~\ref{lem:fibrewise-iso}.
\end{strategy}

\step{1} Define
\[
\varphi\colon E(\eta)\longrightarrow E\big(\bar g^{\,*}\xi\big),
\qquad
\varphi(e)=\big(\pi_\eta(e),\,g(e)\big).
\]
This is well defined: by definition of $\bar g$ we have
$\pi_\xi(g(e))=\bar g(\pi_\eta(e))$, which is exactly the condition defining the
pull-back. It is continuous because $\pi_\eta$ and $g$ are, and it commutes with
the projections to $B(\eta)$.

\step{2} On the fibre over $b$, the map $\varphi$ is
$e\mapsto(b,g(e))$, and $g$ restricts to an isomorphism
$F_b(\eta)\to F_{\bar g(b)}(\xi)$, which is by definition the fibre of
$\bar g^{\,*}\xi$ over $b$.

\step{3} By Lemma~\ref{lem:fibrewise-iso}, $\varphi$ is an isomorphism of
bundles.
\end{proof}

\begin{construction}[Products and Whitney sums]
For $\xi_1,\xi_2$ over $B_1,B_2$ the \emph{Cartesian product}
$\xi_1\times\xi_2$ over $B_1\times B_2$ has total space
$E(\xi_1)\times E(\xi_2)$, projection $\pi_1\times\pi_2$ and fibres
$F_{b_1}(\xi_1)\times F_{b_2}(\xi_2)$; products of local trivializations
trivialize it. When $B_1=B_2=B$, pulling back along the diagonal
$d\colon B\to B\times B$, $b\mapsto(b,b)$, gives the \emph{Whitney sum}
\[
\xi_1\oplus\xi_2:=d^*(\xi_1\times\xi_2),
\qquad
F_b(\xi_1\oplus\xi_2)\iso F_b(\xi_1)\oplus F_b(\xi_2).
\]
\end{construction}

\begin{definition}
Let $\xi,\eta$ be bundles over $B$ with $E(\xi)\subset E(\eta)$. We call $\xi$ a
\emph{sub-bundle} of $\eta$ if every $F_b(\xi)$ is a vector subspace of
$F_b(\eta)$.
\end{definition}

\begin{lemma}\label{lem:internal-sum}
If $\xi_1,\xi_2$ are sub-bundles of $\eta$ such that
$F_b(\eta)=F_b(\xi_1)\oplus F_b(\xi_2)$ for every $b$, then
$\eta\iso\xi_1\oplus\xi_2$.
\end{lemma}

\begin{proof}
The map $E(\xi_1\oplus\xi_2)\to E(\eta)$ sending $(b,e_1,e_2)\mapsto e_1+e_2$ is
continuous, commutes with the projections, and on the fibre over $b$ is the
canonical isomorphism $F_b(\xi_1)\oplus F_b(\xi_2)\to F_b(\eta)$ provided by the
hypothesis. Apply Lemma~\ref{lem:fibrewise-iso}.
\end{proof}

\section{Euclidean vector bundles}

\begin{definition}
A \emph{Euclidean vector bundle} is a real vector bundle $\xi$ together with a
continuous function $\mu\colon E(\xi)\to\R$ which is quadratic and positive
definite on every fibre. Polarization turns $\mu$ into an inner product on each
fibre, varying continuously with the base point; we call it a \emph{Euclidean
metric} on $\xi$, or a \emph{Riemannian metric} when $\xi=TM$.
\end{definition}

The space $\R^n$ carries its standard metric, which restricts to any
$M\subset\R^n$. On a trivial bundle with a Euclidean metric, Gram--Schmidt --- a
continuous process --- turns the standard frame into $n$ orthonormal nowhere
vanishing sections.

\begin{definition}
Let $\eta$ be a Euclidean bundle of rank $n$ and $\xi\subset\eta$ a sub-bundle
of rank $k$. The \emph{orthogonal complement} $\xi^\perp\subset\eta$ is the
union over $b$ of the subspaces $F_b(\xi)^\perp\subset F_b(\eta)$.
\end{definition}

\begin{theorem}\label{thm:orthogonal-complement}
$\xi^\perp$ is a vector bundle and $\xi\oplus\xi^\perp\iso\eta$.
\end{theorem}

\begin{proof}
\begin{strategy}
Only local triviality of $\xi^\perp$ is at issue, since the second assertion
then follows from Lemma~\ref{lem:internal-sum}. Near a point, build an
orthonormal frame of $\eta$ whose first $k$ vectors span $\xi$; the remaining
$n-k$ vectors then trivialize $\xi^\perp$. Producing such a frame uses two
standard facts: linear independence is an open condition, and Gram--Schmidt is
continuous.
\end{strategy}

\step{1} \emph{Setting up.} Fix $b\in B$ and choose a neighbourhood $U$ over
which both $\xi$ and $\eta$ are trivial. Using a trivialization of $\xi|_U$ and
Gram--Schmidt, choose sections $s_1,\dots,s_k$ of $\xi|_U$ which are orthonormal
at every point of $U$.

\step{2} \emph{Extending the frame at one point.} The vectors
$s_1(b),\dots,s_k(b)$ are linearly independent in $F_b(\eta)$, so we may complete
them to a basis
\[
s_1(b),\dots,s_k(b),\,v_{k+1},\dots,v_n
\]
of $F_b(\eta)$. Using a trivialization of $\eta|_U$, extend the vectors $v_j$ to
sections $s'_{k+1},\dots,s'_n$ of $\eta|_U$ that are constant in that
trivialization.

\step{3} \emph{Independence persists.} In the trivialization of $\eta|_U$ the
determinant of the $n\times n$ matrix with columns
$s_1,\dots,s_k,s'_{k+1},\dots,s'_n$ is a continuous function of the base point
which is non-zero at $b$. Hence there is a neighbourhood $V\ni b$, $V\subset U$,
on which these $n$ sections are nowhere dependent.

\step{4} \emph{Orthonormalizing.} Apply Gram--Schmidt pointwise on $V$ to the
frame of Step 3. Since the first $k$ vectors are already orthonormal, the
process leaves them unchanged and produces sections
$s_1,\dots,s_k,s_{k+1},\dots,s_n$ which form an orthonormal frame of $\eta|_V$
at every point, depend continuously on the base point, and satisfy
$\spann(s_1,\dots,s_k)=F_{b'}(\xi)$ for $b'\in V$. Consequently
$s_{k+1}(b'),\dots,s_n(b')$ is an orthonormal basis of $F_{b'}(\xi)^\perp$.

\step{5} \emph{Local triviality of $\xi^\perp$.} Define
\[
V\times\R^{n-k}\longrightarrow \xi^\perp\big|_V,
\qquad
(b',x)\longmapsto x_1\,s_{k+1}(b')+\dots+x_{n-k}\,s_n(b').
\]
It is continuous, fibrewise a linear isomorphism onto $F_{b'}(\xi)^\perp$ by
Step 4, and its inverse is $e\mapsto(b',(\langle e,s_{k+j}(b')\rangle)_j)$,
continuous because the metric is. Hence $\xi^\perp$ is a vector bundle of rank
$n-k$.

\step{6} \emph{The sum.} By construction
$F_b(\eta)=F_b(\xi)\oplus F_b(\xi^\perp)$ for every $b$, so
Lemma~\ref{lem:internal-sum} gives $\xi\oplus\xi^\perp\iso\eta$.
\end{proof}

\begin{example}\label{ex:normal-bundle}
If $M\subset N$ are smooth manifolds and $N$ is Riemannian, then $TM$ is a
sub-bundle of $TN|_M$ and the \emph{normal bundle of $M$ in $N$} is
$\nu=(TM)^\perp$, so that
\[
TN|_M\ \iso\ TM\oplus\nu .
\]

\begin{center}
\begin{tikzpicture}[scale=0.95]
\draw[thick] (0,0) .. controls (1.2,0.9) and (2.8,-0.6) .. (4,0.3);
\node at (4.4,0.3) {$M$};
\fill (1.9,0.12) circle (1.5pt);
\draw[-{Stealth}] (1.9,0.12) -- (2.85,-0.14);
\node[below] at (2.9,-0.08) {\small $T_xM$};
\draw[-{Stealth}] (1.9,0.12) -- (2.14,1.0);
\node[right] at (2.12,1.0) {\small $\nu_x$};
\end{tikzpicture}
\end{center}
\end{example}

\begin{definition}
A smooth map $f\colon M\to N$ is an \emph{immersion} if
$Df_x\colon T_xM\to T_{f(x)}N$ is injective for every $x$. The figure eight is
the image of an immersion of $S^1$ in $\R^2$ which is not an embedding.
\end{definition}

\begin{lemma}\label{lem:normal-immersion}
For an immersion $f\colon M\to N$ with $N$ Riemannian there is a Whitney sum
decomposition
\[
f^*TN\ \iso\ TM\oplus\nu_f,
\]
where $\nu_f$ is called the \emph{normal bundle of the immersion} $f$.
\end{lemma}

\begin{proof}
The differential $Df$ realizes $TM$ as a sub-bundle of $f^*TN$: it is fibrewise
injective by the definition of an immersion, and it is continuous, so its image
is a sub-bundle. Pulling back the Riemannian metric of $N$ makes $f^*TN$
Euclidean, and Theorem~\ref{thm:orthogonal-complement} applied to
$TM\subset f^*TN$ gives $f^*TN\iso TM\oplus(TM)^\perp$; set $\nu_f=(TM)^\perp$.
\end{proof}

This decomposition is the engine behind many computations to come: it converts
geometric information --- how $M$ sits inside $N$ --- into bundle-theoretic
information. In the next lecture we manufacture new bundles from old ones
functorially, and meet the first characteristic classes.
